\section{Marco Metodológico}

Este capítulo describe el enfoque metodológico adoptado para el desarrollo y la validación del prototipo WebGL del Holybro X500~V2. La investigación se concibe como aplicada, con enfoque mixto y predominio cuantitativo, y se apoya en Design Science Research (DSR) para justificar la construcción y evaluación del artefacto.

\subsection{Tipo de investigación y enfoque}

\begin{itemize}
    \item \textbf{Tipo de investigación:} aplicada.
    \item \textbf{Enfoque:} mixto con predominio cuantitativo.
    \item \textbf{Marco:} Design Science Research complementado con prototipado iterativo.
    \item \textbf{Unidad de análisis:} la interacción de usuarios con un visor WebGL orientado a la exploración técnica del Holybro X500~V2.
\end{itemize}

\subsection{Fases de trabajo}

\begin{enumerate}
    \item \textbf{Problema y fundamentación.} Revisión de literatura y definición de metas de visualización, rendimiento y experiencia de uso.
    \item \textbf{Pipeline de activos.} Limpieza, optimización e integración del modelo en Unity.
    \item \textbf{Implementación del prototipo.} Desarrollo de UI, selección de piezas, modos de análisis y visualización.
    \item \textbf{Validación y cierre.} Auditorías funcionales, medición de KPIs y evaluación de usabilidad.
\end{enumerate}

\subsection{Muestra y participantes}

Se plantea una muestra por conveniencia de entre 8 y 12 participantes con perfil afín al caso de estudio: estudiantes avanzados, técnicos o profesionales de ingeniería, manufactura digital, mantenimiento o áreas relacionadas. La muestra es coherente con un estudio exploratorio de usabilidad y carga cognitiva.

\subsection{Variables de investigación}

\subsubsection{Variable independiente}

\textbf{Tipo de visualización técnica.}
\begin{itemize}
    \item nivel 1: visualización 3D interactiva WebGL;
    \item nivel 2: soporte 2D de referencia para contraste metodológico.
\end{itemize}

\subsubsection{Variables dependientes}

\begin{itemize}
    \item \textbf{Usabilidad percibida:} puntaje SUS en escala 0--100.
    \item \textbf{Carga cognitiva:} puntaje NASA-TLX Raw en escala 0--100.
    \item \textbf{Rendimiento técnico:} FPS, tiempos de carga, draw calls, memoria y consistencia funcional del build final.
\end{itemize}

\subsubsection{Variables de control}

\begin{itemize}
    \item experiencia previa con visualización 3D o documentación técnica;
    \item perfil académico o profesional;
    \item dispositivo, navegador y condiciones de conexión.
\end{itemize}

\subsection{Instrumentos}
\label{sec:instrumentos}

\begin{itemize}
    \item \textbf{KPIs técnicos:} herramientas de profiling y auditoría sobre la build WebGL.
    \item \textbf{SUS:} System Usability Scale (Brooke, 1996), aplicado como cuestionario posterior a la interacción.
    \item \textbf{NASA-TLX Raw:} medición de carga cognitiva percibida sobre seis dimensiones (Hart \& Staveland, 1988).
    \item \textbf{Think-Aloud:} observación cualitativa concurrente durante la ejecución de tareas.
\end{itemize}

\subsection{Protocolo de tareas}

Las tareas de validación se alinean con el flujo visible final de la aplicación:

\begin{enumerate}
    \item ingresar al visor desde el Hero;
    \item navegar la escena con orbit, pan y zoom;
    \item seleccionar una pieza y revisar su ficha técnica en el \textit{bottom sheet};
    \item activar una herramienta de \texttt{Inspect};
    \item usar \texttt{Explode} o \texttt{Cut} en \texttt{Analyze};
    \item cambiar el modo visual a \texttt{X-Ray} o \texttt{Thermal} en \texttt{Studio}.
\end{enumerate}

\subsection{Procedimiento de análisis}

\subsubsection{Análisis cuantitativo}

El puntaje SUS se calcula mediante el procedimiento estándar de Brooke. Para los ítems impares, la contribución de cada respuesta se obtiene restando 1 al valor marcado por el participante; para los ítems pares, la contribución se obtiene restando la respuesta a 5. La suma de las diez contribuciones produce un subtotal entre 0 y 40, que luego se multiplica por 2.5 para obtener un puntaje global entre 0 y 100:

\[
SUS = \left(\sum_{j \in impares}(R_j - 1) + \sum_{j \in pares}(5 - R_j)\right)\times 2.5
\]

donde $R_j$ es la respuesta del participante en el ítem $j$. Un valor alto indica mejor usabilidad percibida. Como referencia interpretativa, un puntaje de 68 suele considerarse el umbral convencional de aceptabilidad general, aunque la lectura final debe complementarse con el contexto del caso de uso y con escalas adjetivales como las discutidas por Bangor et al. (2009).

El puntaje global de NASA-TLX Raw se obtiene mediante:
\[
NASA\text{-}TLX_{Raw} = \frac{1}{6}\sum_{i=1}^{6} D_i
\]

donde $D_i$ representa el puntaje asignado a la dimensión $i$ del instrumento: demanda mental, demanda física, demanda temporal, rendimiento, esfuerzo y frustración. En esta adaptación se utiliza la variante \textit{Raw TLX}, por lo que el resultado corresponde al promedio aritmético simple de las seis escalas. La dimensión de rendimiento se diligencia con orientación invertida (0 = perfecto, 100 = fallido), de modo que un valor más alto siga indicando mayor carga de trabajo percibida y pueda promediarse directamente con las demás dimensiones.

Los KPIs técnicos se contrastan con las metas operativas fijadas para el proyecto, sin cerrar resultados hasta disponer del build congelado posterior a la optimización final.

\subsubsection{Análisis cualitativo}

El protocolo \textit{Think-Aloud} se aplica en modalidad concurrente durante la ejecución de tareas. A cada participante se le solicita verbalizar qué intenta hacer, cómo interpreta la interfaz, qué dudas le surgen, qué errores percibe y con qué criterio decide su siguiente acción. El evaluador registra verbalizaciones relevantes, pausas, puntos de vacilación, errores de navegación, solicitudes de ayuda y estrategias de recuperación.

Posteriormente, las observaciones se codifican por categorías analíticas: comprensión espacial del ensamblaje, navegación y control, interpretación de la interfaz, correspondencia terminológica, errores o bloqueos y sugerencias de mejora. Esta lectura cualitativa permite triangular los resultados de SUS y NASA-TLX y explicar por qué determinadas tareas resultaron más claras, más confusas o más demandantes para el usuario (Ericsson \& Simon, 1993).

\subsection{Consideraciones éticas}

La participación en pruebas es voluntaria. No se recolectan datos sensibles y los resultados se reportan de forma agregada. El estudio no contempla intervenciones físicas ni exposición a riesgos adicionales para los participantes.

\newpage
